\documentclass[11pt]{article}
\usepackage{enumitem}
\usepackage{xcolor}
\usepackage{fancyhdr}

%configure options for the hyperref package,
%which is used to add hyperlinks to documents
\usepackage{hyperref}
\hypersetup{
    colorlinks=true,
    linkcolor=blue,
    filecolor=magenta,
    urlcolor=blue
    }
\urlstyle{same}

%configure options for the page size
\usepackage[
%paperheight=16cm, paperwidth=12cm,% Set the height and width of the paper
includehead,
nomarginpar,% We don't want any margin paragraphs
%textwidth=10cm,% Set \textwidth to 10cm
%headheight=10mm,% Set \headheight to 10mm
]{geometry}


% Set the page style to "fancy"...
\pagestyle{fancy}
%... then configure it.
\fancyhead{} % clear all header fields
\fancyhead[L]{TKRR25 Project Report,  HT2025} % left header
\fancyhead[C]{} % center header
\fancyhead[R]{\thepage} % right header

\title{Sudoku Game with ASP-based Hint System}
\author{Emil Fröding, Erik Wöhry, George Saba, Rebecca Huynh}
%\date{}
\begin{document}

\maketitle  % This command generates the title
\newpage
\section{Introduction and Movivation}

 note: Clearly define the scope of your project (what you cover, what you do not). inte klar, första layout

The aim of this project is to demonstrate how Answer Set Programming can be used to analyze puzzle states, detect unsolvable solution/configurations?, and generate hints during gameplay. We will implement an interactive version of Chain Sudoku, also known as Strimko, supported by an ASP-based reasoning component. In Chain sudoku, the traditional 3x3 boxes are replaced by nine irregular "chains", each of which must contain the digits 1-9 exactly once. This gives us a great opportunity to adapt the standard Sudoku encoding and show how ASP can flexibly capture alternative constraint structures. This project will also feature a simple React frontend to accomplish the interactive part. By forwarding the player's current entries to the ASP solver. The frontend is intentionally minimal: it does not validate each move, enforce Sudoku rules during input, or allow penciling notes, as the focus is on the reasoning capabilities rather than interface complexity. The solver will evaluate whether the puzzle remains solvable based on the player's inputs and can provide hints by computing a valid completed solution. 

\newpage
\section{Methodology}

Content depends on your chosen track, but structure is similar

\subsection*{ASP}

{\it 
\begin{itemize}
  \item Describe the ASP encoding: main predicates, rules, constraints.
  \item Explain how the ASP part interacts with the ontology/KG (if applicable).
  \item Explain why ASP is suited for this reasoning.
\end{itemize} 
}
\newpage
\section{Evaluation}

{\it Describe how you evaluated your system (according to the evaluation suggestions given in the project description), and present the evaluation results.}

\newpage
\section{Discussion and Conclusion}

{\it
\begin{itemize}
  \item Discuss the strengths and limitations of your approach.
  \item What did you learn about combining symbolic and data-driven AI?
\end{itemize}
}

\newpage
\begin{thebibliography}{99}

\bibitem{}


\end{thebibliography}



\end{document}
